% Options for packages loaded elsewhere
\PassOptionsToPackage{unicode}{hyperref}
\PassOptionsToPackage{hyphens}{url}
%
\documentclass[
]{article}
\usepackage{amsmath,amssymb}
\usepackage{lmodern}
\usepackage{iftex}
\ifPDFTeX
  \usepackage[T1]{fontenc}
  \usepackage[utf8]{inputenc}
  \usepackage{textcomp} % provide euro and other symbols
\else % if luatex or xetex
  \usepackage{unicode-math}
  \defaultfontfeatures{Scale=MatchLowercase}
  \defaultfontfeatures[\rmfamily]{Ligatures=TeX,Scale=1}
\fi
% Use upquote if available, for straight quotes in verbatim environments
\IfFileExists{upquote.sty}{\usepackage{upquote}}{}
\IfFileExists{microtype.sty}{% use microtype if available
  \usepackage[]{microtype}
  \UseMicrotypeSet[protrusion]{basicmath} % disable protrusion for tt fonts
}{}
\makeatletter
\@ifundefined{KOMAClassName}{% if non-KOMA class
  \IfFileExists{parskip.sty}{%
    \usepackage{parskip}
  }{% else
    \setlength{\parindent}{0pt}
    \setlength{\parskip}{6pt plus 2pt minus 1pt}}
}{% if KOMA class
  \KOMAoptions{parskip=half}}
\makeatother
\usepackage{xcolor}
\setlength{\emergencystretch}{3em} % prevent overfull lines
\providecommand{\tightlist}{%
  \setlength{\itemsep}{0pt}\setlength{\parskip}{0pt}}
\setcounter{secnumdepth}{-\maxdimen} % remove section numbering
\ifLuaTeX
  \usepackage{selnolig}  % disable illegal ligatures
\fi
\IfFileExists{bookmark.sty}{\usepackage{bookmark}}{\usepackage{hyperref}}
\IfFileExists{xurl.sty}{\usepackage{xurl}}{} % add URL line breaks if available
\urlstyle{same} % disable monospaced font for URLs
\hypersetup{
  hidelinks,
  pdfcreator={LaTeX via pandoc}}

\author{Dietmar Gerald Schrausser}
\date{09.05.2026}


\begin{document}

\hypertarget{foliumblat-iiiv}{%
\section{Foliu{[}m{]}/Blat IIIv}\label{foliumblat-iiiv}}

Eine zeitgenössische deutsche Übersetzung wird präsentiert, die auf
einem Vergleich der bestehenden lateinischen, frühneuhochdeutschen
(ENHG) und englischen (Übersetzungs-) Texte basiert, um ein besseres
Verständnis zu ermöglichen. Insbesondere da der ENHG-Text im Vergleich
zum modernen Hochdeutschen schwer verständlich ist, da dieser fundierte
Deutschkenntnisse (als Muttersprache) sowie Kenntnisse verschiedener
Dialekte voraussetzt.

\hypertarget{vom-werk-des-dritten-tages}{%
\subsection{Vom Werk des dritten
Tages}\label{vom-werk-des-dritten-tages}}

\begin{quote}
\textbf{T}Ercio die deus aquas sub firmamento in locum vnum
co{[}n{]}gregauit {[}et{]} appareat arida: vocauitq{[}ue{]} aridam
terram. Congretat{[}i{]}ones vero aquarum maria appellauit. Videns
q{[}uia{]} bonum esset: ait germinet terra herbam virentem: {[}et{]}
semen facientem: {[}et{]} lignum pomiferum facientem fructum iuxta
gen{[}us{]} suum.
\end{quote}

Nach dem Firmament zur Integrität, Position und Ordnung der Elemente,
erinnert er uns kurz an das Zusammenfließen der Wasser an \emph{einen}
Ort und die Gesetze, welche dem Meer vorgeschrieben, so es nicht das
Land überflute.\footnote{vom `Heptaplus' (Mirandola \& Mirandola,
  \href{https://doi.org/10.3931/e-rara-61217}{1601}, Buch 3).}

War die Erde einst unsichtbar und soll diese in das
\emph{Angesicht}\footnote{`conspectum' und `gesycht'.} kommen, so ist es
notwendig, dass die Wasser, die sich unter dem Himmel bzw. unter der
mittleren Luftschicht befinden, an einen Ort \emph{zusammengeführt}
werden, in ein gemeinsames Mischbecken, wo sie gemäß bestimmter Gesetze
zusammenfließen, wie auch für die Küsten quasi vorgeschrieben ist, sich
an einer Stelle zu einem \emph{Damm}\footnote{`matricem'.}
zusammenzufügen.\footnote{vom `Heptaplus' (Mirandola \& Mirandola,
  \href{https://doi.org/10.3931/e-rara-61217}{1601}, Buch 2, Kapitel 5 /
  4).}

Unwahr ist nämlich, dass nirgendwo an abgelegenen und isolierten Orten
Wasser gefunden wird, da das \emph{indische} Meer vom hircanischen Meer,
das \emph{hircanische} Meer vom adriatischen Meer, das
\emph{adriatische} Meer vom \emph{euxinischen} Meer, neben unzähligen
entlegen fließenden \textbf{\emph{Quelle}} \textbf{\emph{n un}}
\textbf{\emph{d Seen}}, durch Wehre\footnote{`intercapedine' und
  `verre'.} von einander geschieden wird. Nachdem es sich hierbei um
eine gesonderte und unterteilte Ansammlung von Meeres- oder
Flussgewässern handelt, \emph{sollten sie wohl an einem gemeinsamen Ort
versammelt sein}\footnote{`Sed \emph{ideo} ad locum vnum congregate aque
  dicu{[}n{]}tur' und `Aber \emph{darumb} werden die wasser an ainige
  stat versamelt genant'.}.\footnote{vom `Heptaplus' (Mirandola \&
  Mirandola, \href{https://doi.org/10.3931/e-rara-61217}{1601}, Buch 2,
  Kapitel 2).} \emph{Alle}, wie \textbf{Salomon} sagt, tendieren zum
primären Meer, werden an einen Ort der Ozeane geführt und
vereint.\footnote{Bezug auf König Salomons \emph{höchst bedeutsamen}
  Vers in Prediger 1,7, `Omnia flumina intrant in mare, et mare non
  redundat; ad locum unde exeunt flumina revertuntur' und weiter `ut
  iterum fluant' (vgl. Jiménez de Cisneros,
  \href{https://doi.org/10.3931/e-rara-46695}{1517}). `Alle Flüsse
  münden ins Meer und doch wird das Meer nie voll, zu dem Ort, von wo
  die Flüsse herkommen, kehren sie zurück - um {[}immer{]} wieder zu
  fließen'. \emph{Das} beschreibt die endlose, zyklische Natur der Welt.}

Wird die Erde\footnote{`terra' und `erde'.} dann von den Wellen des
Meeres überwältigt, ist diese vorerst nicht nützlich oder brauchbar für
uns, wird jedoch für Tiere und auch für unsere Zwecke geeignet, wenn das
Meer \emph{zurückweicht}, dann wird die Fruchtbarkeit und
Ertragreichheit immer klarer ersichtlich. Was auch von \textbf{Moses}
sehr stark verdeutlicht wird, wenn er sie als
\emph{Elternteil}\footnote{`parentem' und `gepererein'.} von Kräutern,
Früchten und Bäumen darstellt und diese nach der Wasseransammlung als
sogleich grün und blühend beschreibt.\footnote{deutsche Renaissance
  Botaniker, wie Otto Brunfels
  (\href{https://www.digitale-sammlungen.de/de/view/bsb11200240?page=1}{1532})
  in seinem bahnbrechenden Werk `Herbarum vivae eicones', zitieren diese
  Stelle aus dem `Heptaplus' (Mirandola \& Mirandola,
  \href{https://doi.org/10.3931/e-rara-61217}{1601}, Buch 2, Kapitel 4)
  häufig um ihre Pflanzenstudien theologisch und philosophisch zu
  untermauern.}

So platzierte er sie (die \emph{Erde}) in die Mitte der \emph{Welt}, als
das \emph{Zentrum}, mit Adern\footnote{`venis'.} aus Metallen wie Gold,
Silber, Bronze\footnote{`ere', aeris.}, Kupfer, Zinn, Blei und (alle
übertreffend) Eisen ausgestattet, sofort mit allen Arten von Kräutern,
mit größter Freude in ihrer grünen Reife, bedeckt, mit Kräutersamen und
Bäume mit den süßesten Früchten. Es heißt auch, dass er an diesem Tag
den fruchtbarsten und herrlichsten \emph{Ursprung}\footnote{`ortum' und
  `garten'.}, das Paradies erschaffen habe, mit allen Arten von Hölzern
und Bäumen, zusammengestellt aus allen Annehmlichkeiten von Quellen, mit
grünen Landstrichen von Bäumen, die reichlichsten Früchte
tragend.\footnote{ein Auszug aus dem `Supplementum Chronicarum' (s.
  Foresti,
  \href{https://books.google.com/books?id=ei9TruMbYCkC\&printsec=frontcover}{1492}).}

\hypertarget{references}{%
\subsection{References}\label{references}}

Brunfels, O. (1532). \emph{Herbarum Vivae Eicones: Ad Naturae
Imitationem Summa Cum Diligentia \& Artificio Effigiatae, Una Cum
Effectibus Earundem \ldots{} ; Quibus Adiecta Ad Calcem, Appendix
Isagogica de Usu \& Administratione Simplicium}. Argentorati: apud
Ioannem Schottum.
\url{https://www.digitale-sammlungen.de/de/view/bsb11200240?page=1}

Foresti, G. F. (1492). \emph{Supplementum Chronicarum}. Novariensis:
Bernardinus Rizus.
\url{https://books.google.com/books?id=ei9TruMbYCkC\&printsec=frontcover}

Jiménez de Cisneros, F. (1517). \emph{Biblia Polyglotta Complutensis}.
Complutum: Arnaldo Guillén de Brocar.
\url{https://doi.org/10.3931/e-rara-46695}

Mirandola, G., \& Mirandola, G. F. (1601). \emph{Ioannis Pici,
Mirandulae \ldots{} opera quae extant omnia \ldots{}} Basileae: per
Sebastianum Henricpetri. \url{https://doi.org/10.3931/e-rara-61217}

\end{document}
